\section{Interferometria Neutronica}
In questo esperimento si invia un fascio di neutroni attraverso un beam splitter che li divide in due fasci, uno libero e l'altro passante attraverso un solenoide. I due fasci vengono poi ricombinati e si osserva il pattern di interferenza risultante. L'obiettivo è la verifica dell'andamento degli spinori (dove una rotazione di 2$\pi$ equivale a un cambiamento di segno). L'hamiltoniana risulta essere
\[
\H = \omega \SG z;\ \ \omega = \frac{eB}{m_n C}g_n
\]
dove $g_n$ è un fattore di correzione.

Lo stato di partenza è $\Ket +$, e lo stato dei fasci risulta essere
\[
\Ket {\psi_1(t)} = \Ket {\psi_1(0)} = e^{i\delta}\Ket +\]\[
\Ket {\psi_2(t)} = e^{-i \frac{\omega t}{2}} \Ket + \]\[
\Ket {\psi_s(t)} = \half (\Ket {\psi_1} + \Ket {\psi_1}) = \half \left( e^{i\delta} + e^{-i \frac{\omega t}{2}} \right) \Ket + + \half \left( e^{i\delta} + e^{i \frac{\omega t}{2}} \right) \Ket -
\]
$\delta$ è una differenza di fase indotta dal beam splitter. E risulta (calcoli per esercizio):
\[
    P_+ = |\braket{+}{\psi_s(t)}|^2 = \half\left(1 - \cos\left( \frac{\omega t}{2} + \delta \right)\right)
\]

% \Ket{\psi(t)} = \frac{1}{\sqrt{2}}\left( \Ket{+} + e^{\frac{-i}{\hslash}\omega t \SG z}\Ket{+} \right) = \frac{1}{\sqrt{2}}\left( \Ket{+} + \left( \cos(\omega t) \Id - i\sin(\omega t) \SG z \right)\Ket{+} \right)

\section{Perturbation Theory}
Cose dal quaderno, modalità di perturbation theory.

Sia l'hamiltoniana quindi indipendente dal tempo, a spettro discreto, non degenere (con autostati di moltemplicità 1). Si supponga di alterare l'hamiltoniana $\H$ con un termine di potenziale piccolo $\varepsilon V$. Quel che risulta è che gli autovalori e autostati della nuova hamiltoniana $\H_\varepsilon = \H + \varepsilon V$ si discostano di poco da quelli di $\H$, e in ogni caso in maniera analitica. Si esprime la correzione al primo termine per i nuovi autovalori e i nuovi autostati:
\[
E_n^\varepsilon = E_n^{(0)} + \varepsilon E_n^{(1)} + \varepsilon^2 E_n^{(2)} + \ldots \]\[
\Ket u_n^\varepsilon = \Ket u_n^{(0)} + \varepsilon \Ket u_n^{(1)} + \varepsilon^2 \Ket u_n^{(2)} + \ldots
\]
E quindi
\[
    \H_\varepsilon \Ket u_n^\varepsilon = E_n^\varepsilon \Ket u_n^\varepsilon
\]
Se ci si accontenta dell'espansione all'ordine lineare risulta:
\[
    (\H + \varepsilon V)(\Ket u_n^{(0)} + \varepsilon \Ket u_n^{(1)}) = (E_n^{(0)} + \varepsilon E_n^{(1)})(\Ket u_n^{(0)} + \varepsilon \Ket u_n^{(1)}) = E_n^{(0)} \Ket u_n^{(0)} + \varepsilon \left( E_n^{(0)} \Ket u_n^{(1)} + E_n^{(1)} \Ket u_n^{(0)} \right)
\]
Oppure, rimuovendo anche i termini di ordine zero e dividendo per $\varepsilon$:
\[
    \H \Ket u_n^{(1)} + V \Ket u_n^{(0)} = E_n^{(0)} \Ket u_n^{(1)} + E_n^{(1)} \Ket u_n^{(0)}\]\[
    E_n^{(1)} = \bra u_n^{(0)}| V |u_n^{(0)}\ket
\]
che è un ottimo punto di partenza per calcolare la prima correzione degli autovalori.
Ora si può imporre che la proiezione del primo stato variato sullo stato invariato sia 1: $\bra u_n^{(0)}|u_n^\varepsilon\ket = 1$, cosicché si ottiene che gli stati successivi al primo sono ortogonali al primo: $\bra u_n^{(0)}|u_n^{(1)}\ket = 0$ (verificare espandendo l'espressione precedente).

Si può procedere in maniera ricorsiva per ottenere gli stati variati di ordine successivo, a partire dal secondo:
\[
    \H \Ket u_n^{(2)} + V \Ket u_n^{(1)} = E_n^{(0)} \Ket {u_n^{(2)}} + E_n^{(1)} \Ket {u_n^{(1)}} + E_n^{(2)} \Ket {u_n^{(0)}}
\]

E di nuovo, completando il prodotto scalare a sinistra con $\bra u_n^{(0)}|$ si ottiene:
\[
    \bra u_n^{(0)}| V |u_n^{(1)}\ket = E_n^{(0)} \bra u_n^{(0)}| u_n^{(2)}\ket + E_n^{(1)} \bra u_n^{(0)}| u_n^{(1)}\ket + E_n^{(2)} \bra u_n^{(0)}| u_n^{(0)}\ket = E_n^{(2)}
\]

Ora si moltiplichi questo risultato per $\bra u_m^{(0)}|$ con $m \neq n$:
\[
    \bra u_m^{(0)}| \H \Ket {u_n^{(2)}} + \Bra {u_m^{(0)}} V \Ket u_n^{(1)} = E_n^{(0)} \bra u_m^{(0)}| u_n^{(2)}\ket + E_n^{(1)} \bra u_m^{(0)}| u_n^{(1)}\ket + E_n^{(2)} \bra u_m^{(0)}| u_n^{(0)}\ket
\]
E sapendo che $\H \Ket {u_m^{(0)}} = E_m^{(0)} \Ket {u_m^{(0)}}$ e che $\bra u_m^{(0)}| u_n^{(0)}\ket = 0$ si ottiene:
\[
    (E_m^{(0)} - E_n^{(0)}) \bra u_m^{(0)}| u_n^{(1)}\ket = - \Bra {u_m^{(0)}} V \Ket u_n^{(0)}
\]

E quindi abbiamo ottenuto la proiezione della prima correzione dello stato $n$ sullo stato generico della base degli stati invariati:
\[
    \bra u_m^{(0)}| u_n^{(1)}\ket = \sum_{m \neq n} \frac{|\Bra {u_m^{(0)}} V \Ket u_n^{(0)}|^2}{E_n^{(0)} - E_m^{(0)}} \ \implies\ \Ket {u_n^{(1)}} = \sum_{m \neq n} \frac{|\Bra {u_m^{(0)}} V \Ket u_n^{(0)}|^2}{E_n^{(0)} - E_m^{(0)}} \Ket {u_m^{(0)}}
\]
Al ché è come se $\Bra {u_m^{(0)}} V \Ket u_n^{(0)}$ fosse una matrice degli elementi di perturbazione tra gli stati $m$ ed $n$.
Gli elementi di questa matrice rappresentano, osservando la natura del prodotto scalare, la sovrapposizione tra lo i due stati considerati, ovvero l'influenza dello stato $m$ sullo stato $n$ attraverso la perturbazione $V$, pensando con la differenza di energia dei livelli.

Quel che otteniamo è, all'atto pratico, una quantificazione dell'influenza di una perturbazione esterna sugli stati.

Ora possiamo assemblare la ricetta per esprimere $E_n^{(2)}$:
\[
    E_n^{(2)} = \bra u_n^{(0)}| V |u_n^{(1)}\ket = \bra u_n^{(0)}| V |\left(\sum_{m \neq n} \frac{|\Bra {u_m^{(0)}} V \Ket u_n^{(0)}|^2}{E_n^{(0)} - E_m^{(0)}}\right) \Ket {u_n^{(0)}} = \sum_{m \neq n} \frac{|\Bra {u_m^{(0)}} V \Ket u_n^{(0)}|^2}{E_n^{(0)} - E_m^{(0)}}
\]
-
Risulta poi che $\Ket {u_n^{(2)}}$ può essere espresso come:
\[
    \Ket {u_n^{(2)}} = \sum_{m \neq n} \frac{\Bra {u_m^{(0)}} V \Ket {u_n^{(1)}}}{E_n^{(0)} - E_m^{(0)}} \Ket {u_m^{(0)}} = \sum_{m \neq n} \left( \sum_{k \neq n} \frac{\Bra {u_k^{(0)}} V \Ket {u_n^{(0)}} \Bra {u_m^{(0)}} V \Ket {u_k^{(0)}}}{(E_n^{(0)} - E_k^{(0)})(E_n^{(0)} - E_m^{(0)})} \right) \Ket {u_m^{(0)}}
\]

Quanto vale la norma di $\Ket {u_n^\epsilon}$?
\[
    ||\Ket {u_n^\epsilon}||^2 = \braket {u_n^\epsilon}{u_n^\epsilon} = \left( \Bra {u_n^{(0)}} + \varepsilon \Bra {u_n^{(1)}} \right) \left( \Ket {u_n^{(0)}} + \varepsilon \Ket {u_n^{(1)}} \right) = 1 + \varepsilon^2 \braket {u_n^{(1)}}{u_n^{(1)}}
\]
Che va necessariamente normalizzato in questo modo:
\[
    \Ket {\tilde{u}_n^\epsilon} = \frac{\Ket {u_n^\epsilon}}{||\Ket {u_n^\epsilon}||} = \left( 1 + \varepsilon^2 \braket {u_n^{(1)}}{u_n^{(1)}} \right)^{-\half} \left( \Ket {u_n^\epsilon} \right) \simeq \left( 1 - \half \varepsilon^2 \braket {u_n^{(1)}}{u_n^{(1)}} \right) \left( \Ket {u_n^\epsilon} \right)
\]
Il termine di correzione alla norma è chiamato \textbf{Loss Term} \index{Loss Term}.